\documentclass[12pt]{article}
\usepackage[utf8]{inputenc}
\usepackage[T1]{fontenc}
\usepackage{hyperref}
\title{Kirya-Konzel Language Resources for K-12 Curriculum Development}
\author{Compiled by Claude Sonnet 5.5 (claude-sonnet-5-5) for Andr\'as Kornai}
\date{}
\begin{document}
\maketitle

\section{Introduction}
Kirya-Konzel (ISO 639-3 fkk) is a Chadic (Biu--Mandara, Higi) language of Adamawa State, Nigeria, with varieties Kirya and Konzel that are very close. Wikipedia gives about 5,800 native speakers (2011); Joshua Project lists 8,600 people. Hausa is the language of wider communication. Little exists; this report is based on a single search.

\section{Primary Educational Resources}
None found.

\section{Linguistic and Reference Materials}
\begin{itemize}
\item Wikipedia, Kirya-Konzel language: \url{https://en.wikipedia.org/wiki/Kirya-Konz%C9%99l_language}
\item Wikitongues language page: \url{https://wikitongues.org/languages/fkk/}
\end{itemize}

\section{Religious and Community Resources}
Joshua Project reports only Bible portions in the language; many speakers read Hausa:
\url{https://joshuaproject.net/people_groups/20568} and \url{https://joshuaproject.net/index.php/languages/fkk}

\section{Digital and Audio Resources}
None found.

\section{Curriculum Development Context}
No textbooks, primers or school programmes were found. Literacy is probably in Hausa. A curriculum would need new orthography and materials work with the community.

\end{document}